% !TEX program = xelatex
% 中文双栏论文模板（模仿 IEEEtran 风格）
% Overleaf 使用：New Project → 上传本文件，菜单中编译器选 XeLaTeX 即可
% 字体：ctex 自动检测环境（Overleaf 用自带 Fandol 字体，无需配置）
\documentclass[10pt,twocolumn]{ctexart}

\usepackage[a4paper,top=1.9cm,bottom=2.2cm,left=1.57cm,right=1.57cm,columnsep=0.6cm]{geometry}
\usepackage{amsmath,amsfonts,amssymb}
\usepackage{graphicx}
\usepackage{booktabs}
\usepackage{multirow}
\usepackage{newtxmath}             % 数学字体用 newtx（Times 风格）
\setmainfont{TeX Gyre Termes}      % 正文西文用 TeX Gyre Termes（Times 系）
% 注意：不要用 newtxtext，它会重定义 \ttfamily 导致 xeCJK 中文等宽字体（仿宋）加载失败
\usepackage{caption}
\usepackage[colorlinks=false,pdfborder={0 0 0}]{hyperref}

% ---------- 页面与版式 ----------
\pagestyle{plain}
\setlength{\parindent}{1.5em}
\linespread{1.05}
\captionsetup{font=small,labelfont=bf}

% ---------- 标题样式（模仿 IEEE：大标题居中、作者小字） ----------
\usepackage{titling}
\pretitle{\begin{center}\LARGE\bfseries}
\posttitle{\par\end{center}\vskip 0.6em}
\preauthor{\begin{center}\normalsize}
\postauthor{\par\end{center}\vskip 0.5em}
\predate{\begin{center}\small}
\postdate{\par\end{center}\vskip 1em}

% ---------- 章节标题：模仿 IEEE 的罗马数字大写 ----------
\ctexset{
  section = {
    name = {,.},           % 输出 "I." "II." 形式
    number = \Roman{section},
    format = \centering\scshape,   % IEEE 惯例：居中、小型大写（中文退化为宋体加间距）
    aftername = {\ },
    beforeskip = 1.2em plus 0.2em,
    afterskip = 0.8em plus 0.1em,
  },
  subsection = {
    name = {,.},
    number = \Alph{subsection},   % IEEE 惯例：子节用字母编号 A. B. C.
    format = \itshape\raggedright,
    aftername = {\ },
    beforeskip = 1.0em plus 0.2em,
    afterskip = 0.5em plus 0.1em,
  },
}

% ---------- 数学宏 ----------
\DeclareMathOperator*{\argmax}{arg\,max}
\DeclareMathOperator{\NMF}{NMF}
% 记号宏（与 method_full.tex 导言区一致，供方法节各分文件使用）
\newcommand{\E}{\mathbb{E}}
\newcommand{\Var}{\operatorname{Var}}
\newcommand{\Cov}{\operatorname{Cov}}
\newcommand{\Bias}{\operatorname{Bias}}
\newcommand{\Pbb}{\mathbb{P}}
\newcommand{\N}{\mathcal{N}}
\newcommand{\BetaDist}{\operatorname{Beta}}
\newcommand{\logit}{\operatorname{logit}}

% ---------- 摘要与关键词（IEEE 风格：加粗"摘要—"+斜体正文） ----------
\newcommand{\cnabstract}[1]{%
  \begin{center}\begin{minipage}{0.92\linewidth}
    \small\itshape
    \noindent\textbf{摘要——}#1
  \end{minipage}\end{center}\vskip 0.5em}
\newcommand{\cnkeywords}[1]{%
  \begin{center}\begin{minipage}{0.92\linewidth}
    \small\itshape
    \noindent\textbf{关键词——}#1
  \end{minipage}\end{center}\vskip 1em}

\begin{document}

% ================= 标题区 =================
\twocolumn[{%
  \centering
  {\LARGE\bfseries 一份模仿 IEEE 风格的中文双栏论文模板\par}
  \vskip 0.8em
  {\normalsize 张~三$^{1}$，李~四$^{2}$\par}
  {\small $^{1}$某某大学 电子工程学院，北京，中国\quad $^{2}$某某研究所，上海，中国\par}
  \vskip 1.5em
}]

% ================= 摘要 =================
\cnabstract{本模板用 \LaTeX\ 的 ctexart 文档类实现中文双栏排版，版式要素模仿 IEEE Transactions：双栏正文、罗马数字章节号、居中小型大写节标题、加粗“摘要——”引导的斜体摘要，以及 IEEE 风格的参考文献编号。本段文字用于展示摘要的排版效果，可替换为论文的真实摘要。}

\cnkeywords{中文排版；双栏模板；\LaTeX；IEEEtran 风格}

% ================= 正文 =================
\section{引言}
这里是引言部分。IEEE 论文通常以一句话背景开场，随后指出研究空白，最后给出本文贡献。中文写作时建议保持同样的叙事结构：背景——问题——贡献列表。

IEEEtran 里的 \texttt{\textbackslash IEEEPARstart} 首字下沉效果是英文专属技巧，中文模板中不使用，直接正常段落起排即可。

% ================= 方法章节 =================
% 由 main.tex 抽出，通过 % ================= 方法章节 =================
% 由 main.tex 抽出，通过 % ================= 方法章节 =================
% 由 main.tex 抽出，通过 \input{method} 引入。
% 依赖 main.tex 导言区：\DeclareMathOperator*{\argmax}{arg\,max} 与 \DeclareMathOperator{\NMF}{NMF}

\section{方法}
\label{sec:method}

\subsection{含背景光的直方图观测模型}
\label{sec:spad}

设 $p$ 为目标像素，$q$ 为其邻域 $\mathcal{N}_p$ 中的邻居，$b$ 为直方图的离散 bin 坐标。目标回波与背景光子可视为独立泊松到达过程，故像素 $p$ 在 bin $b$ 的计数满足
\begin{equation}
  H_p(b)\sim\operatorname{Poisson}\!\big(\lambda_p(b)\big),
  \qquad
  \lambda_p(b)=N_p h(b-\mu_p)+\beta_p,
  \label{eq:poisson_obs}
\end{equation}
其中 $\mu_p$ 为真实回波中心坐标，$N_p$ 为期望信号光子数，$h(\cdot)$ 为归一化系统脉冲响应（IRF），$\beta_p$ 为每 bin 平均背景与暗计数率。将 $\mu_p$ 视为待估参数，含背景泊松模型的 Fisher 信息为
\begin{equation}
  \mathcal{I}_p(\mu_p)
  =N_p^2\sum_b\frac{[h'(b-\mu_p)]^2}{N_p h(b-\mu_p)+\beta_p},
  \label{eq:fisher}
\end{equation}
其中 $h'$ 为 IRF 的导数。以 Fisher 信息的倒数近似回波中心估计方差 $v_p\approx\mathcal{I}_p^{-1}$；当 $\beta_p\rightarrow0$ 且 IRF 为方差 $\sigma_h^2$ 的高斯型时，退化为经典形式
\begin{equation}
  v_p\approx\frac{\sigma_h^2}{N_p},
  \label{eq:v_classical}
\end{equation}
即估计方差与信号光子数成反比；背景光则经式~\eqref{eq:fisher}的分母增大逐 bin 方差、降低估计精度。该反比关系是 \ref{sec:engineering} 节精度因子的构造依据。

\subsection{死时间畸变与 Coates 校正}
\label{sec:coates}

实际 SPAD 在一个测距周期内至多记录最早到达的一个光子，死时间使原始直方图偏离式~\eqref{eq:poisson_obs}的独立泊松模型。设逐 bin 光子到达率为 $r_p(b)$，累积率 $R_p(b)=\sum_{b'\le b}r_p(b')$，帧内有效周期数为 $M$，则原始直方图在 $b$ 处的累积计数满足精确关系 $\mathrm{cum}(b)=M\big(1-e^{-R_p(b)}\big)$，其解析逆（Coates 校正）为
\begin{equation}
  \hat r_p(b)=\ln\frac{M-\mathrm{cum}(b-1)}{M-\mathrm{cum}(b)},
  \label{eq:coates}
\end{equation}
其中 $\mathrm{cum}(b)=\sum_{b'\le b}H_p^{\mathrm{raw}}(b')$ 为原始直方图累积计数。该反演不含待调参数、对任意强度精确可逆：背景平坦底座反演后恢复常数率，窄峰的形状与中心位置近似保持，故校正不参与对齐量的定义，仅改变直方图取值标度与噪声结构。频率代换引入的统计误差为相对量级 $O(1/\sqrt{M})$ 的随机误差，可并入散粒噪声预算；当幸存周期数过少时反演方差失控，本文的测距窗设置使峰区自然满足可信条件。本文后续的对齐、匹配与融合均以校正后的率域直方图 $F_p(b)\triangleq\hat r_p(b)$ 为输入。

\subsection{单目相对深度引导的局部对齐}
\label{sec:alignment}

设 $p$、$q$ 的单目深度坐标为 $b_p^M$、$b_q^M$，定义单目相对位移
\begin{equation}
  \Delta_{pq}\triangleq b_p^M-b_q^M=(\mu_p-\mu_q)+(e_p-e_q),
  \label{eq:delta}
\end{equation}
其中 $e_p$、$e_q$ 为单目预测误差。$\Delta_{pq}$ 承担两个作用：提供邻居直方图局部搜索中心，并作为 $p$、$q$ 是否属于同一表面的几何证据；其中只含误差之差 $e_p-e_q$，单目全局尺度误差在对齐这一步自然抵消。量化到整数 bin 得粗对齐量 $k_{pq}^{M}=\operatorname{round}(\Delta_{pq})$，以其为中心在窗口 $\mathcal{W}=\{-K,\dots,K\}$ 内对候选偏移 $\delta$ 计算归一化匹配滤波得分
\begin{equation}
  \rho_{pq}(\delta)
  =\NMF\!\big(H_p,\;H_q(b-k_{pq}^{M}-\delta)\big),
  \label{eq:nmf}
\end{equation}
其中 $\NMF(\cdot,\cdot)$ 为归一化匹配滤波算子。取 $\hat\delta_{pq}=\argmax_{\delta\in\mathcal{W}}\rho_{pq}(\delta)$，则实际用于对齐的平移量为
\begin{equation}
  \hat k_{pq}=k_{pq}^{M}+\hat\delta_{pq},
  \qquad
  \widetilde H_{pq}(b)=H_q(b-\hat k_{pq}),
  \label{eq:final_shift}
\end{equation}
其中 $\widetilde H_{pq}$ 为对齐到目标坐标系的邻居直方图；最优位置处的归一化得分记为 $\rho_{pq}=\max_\delta\rho_{pq}(\delta)\in(0,1)$。单目先验只负责将搜索限制在合理局部范围内，最终对齐位置由观测数据的匹配峰决定，$\rho_{pq}$ 记录该峰的匹配质量。

\subsection{可融合性后验与贝叶斯融合权重}
\label{sec:fusibility}

定义二值潜变量 $z_{pq}\in\{0,1\}$：$z_{pq}=1$ 表示邻居 $q$ 可融合（对齐后可用于估计 $\mu_p$），$z_{pq}=0$ 表示不可融合（如来自不同表面或噪声巧合峰）。判断证据为 $(\Delta_{pq},\rho_{pq})$。设可融合状态下 $\Delta_{pq}\mid z_{pq}=1\sim\mathcal{N}(0,\sigma_s^2)$，不可融合状态下 $\Delta_{pq}\mid z_{pq}=0\sim\mathcal{N}(0,\sigma_d^2)$ 且 $\sigma_d>\sigma_s$，则几何似然比为
\begin{equation}
  \Lambda_{\Delta}(\Delta_{pq})
  =\frac{\sigma_d}{\sigma_s}
  \exp\!\left[-\frac{\Delta_{pq}^2}{2}\Big(\frac{1}{\sigma_s^2}-\frac{1}{\sigma_d^2}\Big)\right],
  \label{eq:lambda_delta}
\end{equation}
其中 $\sigma_s$、$\sigma_d$ 分别刻画同一表面内单目局部相对误差的集中程度与跨表面时含深度跳变的散布程度。$\Lambda_\Delta$ 随 $|\Delta_{pq}|$ 单调递减，构成连续的几何软门控，而非硬阈值删除。匹配证据经条件 Beta 分布参数化为似然比 $\Lambda_\rho(\rho_{pq}\mid\Delta_{pq})$（有标注数据时可直接拟合）。由 Bayes 定理，后验可融合概率为
\begin{equation}
  \gamma_{pq}
  =\frac{\pi_{pq}\Lambda_{\Delta}\Lambda_{\rho}}
  {(1-\pi_{pq})+\pi_{pq}\Lambda_{\Delta}\Lambda_{\rho}},
  \label{eq:gamma}
\end{equation}
其中 $\pi_{pq}=\mathbb{P}(z_{pq}=1)$ 为先验可融合概率。等价地在 logit 域中两类证据可加：几何提供软门控，匹配验证实际观测支持。

定义最终对齐误差 $\epsilon_{pq}\triangleq\hat k_{pq}-(\mu_p-\mu_q)$，结合式~\eqref{eq:delta}与式~\eqref{eq:final_shift}得三来源分解
\begin{equation}
  \epsilon_{pq}=(e_p-e_q)+r_{pq}+\hat\delta_{pq},
  \label{eq:epsilon}
\end{equation}
其中 $r_{pq}=k_{pq}^{M}-\Delta_{pq}\in[-\tfrac12,\tfrac12]$ 为取整残差。前两项由先验与离散化决定，第三项由观测选出且与前两者一般不独立，故各方差不能直接相加，对齐方差 $\sigma_{\epsilon,pq}^2$ 由配对标注数据或重复观测估计。对齐直方图的回波中心观测为 $y_{pq}=\mu_p+\epsilon_{pq}+n_{pq}$，其中测量噪声 $n_{pq}$ 的方差 $v_q$ 由式~\eqref{eq:fisher}给出。当测量与对齐误差条件独立且对齐偏差已校准时，总有效方差为
\begin{equation}
  v_{pq}^{\mathrm{eff}}=v_q+\sigma_{\epsilon,pq}^{2}.
  \label{eq:veff}
\end{equation}
将可融合状态建模为高斯观测 $y_{pq}\mid z_{pq}=1\sim\mathcal{N}(\mu_p,v_{pq}^{\mathrm{eff}})$，不可融合状态视为与 $\mu_p$ 无关的异常分量；对完整数据负对数似然关于 $z_{pq}$ 取后验期望（利用 $\mathbb{E}[z_{pq}\mid\Delta_{pq},\rho_{pq}]=\gamma_{pq}$），得到邻居的后验期望精度
\begin{equation}
  c_{pq}=\frac{\gamma_{pq}}{v_{pq}^{\mathrm{eff}}}
  =\frac{\gamma_{pq}}{v_q+\sigma_{\epsilon,pq}^{2}},
  \label{eq:c_pq}
\end{equation}
它将``是否值得融合''（$\gamma_{pq}$）与``融合后有多准确''（分母）统一到同一统计量。令归一化系数 $\alpha_{pq}=c_{pq}/\sum_{r}c_{pr}$，融合估计为 $\hat\mu_p=\sum_{q\in\mathcal{N}_p}\alpha_{pq}y_{pq}$。代入式~\eqref{eq:gamma}展开可见，理想权重呈``可靠性（$\gamma_{pq}$）$\times$ 精度（$1/v_{pq}^{\mathrm{eff}}$）''的因子结构，为工程实现给出明确的分解方向。

\subsection{可标定的工程实现}
\label{sec:engineering}

理想权重要求 $\gamma_{pq}$、$v_q$、$\sigma_{\epsilon,pq}^2$ 在运行时逐像素获得，而单次观测中均不可直接观测；可直接计算的只有 $\Delta_{pq}$、$\hat k_{pq}$、$\rho_{pq}$。本文不为其虚构运行时估计，而是构造与理想似然比同形状的代理，全部自由参数由验证数据离线标定。可靠性因子取
\begin{equation}
  g_{\Delta}(\Delta_{pq})=\exp\!\Big(-\frac{\Delta_{pq}^2}{2\sigma_w^2}\Big),
  \qquad
  g_{\rho}(\rho_{pq})=\rho_{\mathrm{floor}}+(1-\rho_{\mathrm{floor}})\rho_{pq},
  \label{eq:reliability}
\end{equation}
其中 $\sigma_w$ 为几何软门控尺度，标定量级对齐单目局部相对误差（对应 $\sigma_s$）；$\rho_{\mathrm{floor}}\in[0,1)$ 为低分邻居的保底贡献，避免匹配得分偶然偏低时硬性删除仍有先验支持的邻居。精度因子保留显式代理
\begin{equation}
  \hat u_q=\frac{1}{\hat N_q+\hat N_0},
  \qquad
  \hat N_q=\sum_{b\in\Omega_q}\big(H_q(b)-\hat\beta_q\big),
  \label{eq:precision_proxy}
\end{equation}
其中 $\hat N_q$ 为直方图峰显著度，即峰值信号窗口 $\Omega_q$ 内扣除背景估计 $\hat\beta_q$ 后的信号光子数代理，由式~\eqref{eq:v_classical}的反比关系对应理想精度因子 $1/v_q$；$\hat N_0>0$ 为正则化项，对应背景光方差下限，防止低信号邻居权重发散。三个标定参数相互耦合较弱，固定搜索窗口半宽 $K$ 与匹配算子后，在验证集上以深度误差（MAE）为目标逐一维搜索确定并冻结。最终融合权重与归一化系数为
\begin{equation}
  w_{pq}=g_{\Delta}(\Delta_{pq})\,g_{\rho}(\rho_{pq})\,\hat u_q,
  \qquad
  \alpha_{pq}=\frac{w_{pq}}{\sum_{r\in\mathcal{N}_p}w_{pr}},
  \label{eq:weight}
\end{equation}
对对齐后的邻域直方图加权融合
\begin{equation}
  H_p^{\mathrm{fuse}}(b)
  =\sum_{q\in\mathcal{N}_p}\alpha_{pq}\,\widetilde H_{pq}(b),
  \label{eq:fuse}
\end{equation}
再由融合直方图提取回波中心，得到最终深度估计 $\hat\mu_p$。当邻域内信号强度变异有限时，可将 $\hat u_q$ 置为常数，退化为仅可靠性加权。整个流程仅使用单目先验与当前 SPAD 观测即可运行，同时保留了理想贝叶斯模型的核心结构：几何一致性抑制跨表面观测，匹配一致性评价对齐可信度，信号强度压低低光子数邻居的贡献。
 引入。
% 依赖 main.tex 导言区：\DeclareMathOperator*{\argmax}{arg\,max} 与 \DeclareMathOperator{\NMF}{NMF}

\section{方法}
\label{sec:method}

\subsection{含背景光的直方图观测模型}
\label{sec:spad}

设 $p$ 为目标像素，$q$ 为其邻域 $\mathcal{N}_p$ 中的邻居，$b$ 为直方图的离散 bin 坐标。目标回波与背景光子可视为独立泊松到达过程，故像素 $p$ 在 bin $b$ 的计数满足
\begin{equation}
  H_p(b)\sim\operatorname{Poisson}\!\big(\lambda_p(b)\big),
  \qquad
  \lambda_p(b)=N_p h(b-\mu_p)+\beta_p,
  \label{eq:poisson_obs}
\end{equation}
其中 $\mu_p$ 为真实回波中心坐标，$N_p$ 为期望信号光子数，$h(\cdot)$ 为归一化系统脉冲响应（IRF），$\beta_p$ 为每 bin 平均背景与暗计数率。将 $\mu_p$ 视为待估参数，含背景泊松模型的 Fisher 信息为
\begin{equation}
  \mathcal{I}_p(\mu_p)
  =N_p^2\sum_b\frac{[h'(b-\mu_p)]^2}{N_p h(b-\mu_p)+\beta_p},
  \label{eq:fisher}
\end{equation}
其中 $h'$ 为 IRF 的导数。以 Fisher 信息的倒数近似回波中心估计方差 $v_p\approx\mathcal{I}_p^{-1}$；当 $\beta_p\rightarrow0$ 且 IRF 为方差 $\sigma_h^2$ 的高斯型时，退化为经典形式
\begin{equation}
  v_p\approx\frac{\sigma_h^2}{N_p},
  \label{eq:v_classical}
\end{equation}
即估计方差与信号光子数成反比；背景光则经式~\eqref{eq:fisher}的分母增大逐 bin 方差、降低估计精度。该反比关系是 \ref{sec:engineering} 节精度因子的构造依据。

\subsection{死时间畸变与 Coates 校正}
\label{sec:coates}

实际 SPAD 在一个测距周期内至多记录最早到达的一个光子，死时间使原始直方图偏离式~\eqref{eq:poisson_obs}的独立泊松模型。设逐 bin 光子到达率为 $r_p(b)$，累积率 $R_p(b)=\sum_{b'\le b}r_p(b')$，帧内有效周期数为 $M$，则原始直方图在 $b$ 处的累积计数满足精确关系 $\mathrm{cum}(b)=M\big(1-e^{-R_p(b)}\big)$，其解析逆（Coates 校正）为
\begin{equation}
  \hat r_p(b)=\ln\frac{M-\mathrm{cum}(b-1)}{M-\mathrm{cum}(b)},
  \label{eq:coates}
\end{equation}
其中 $\mathrm{cum}(b)=\sum_{b'\le b}H_p^{\mathrm{raw}}(b')$ 为原始直方图累积计数。该反演不含待调参数、对任意强度精确可逆：背景平坦底座反演后恢复常数率，窄峰的形状与中心位置近似保持，故校正不参与对齐量的定义，仅改变直方图取值标度与噪声结构。频率代换引入的统计误差为相对量级 $O(1/\sqrt{M})$ 的随机误差，可并入散粒噪声预算；当幸存周期数过少时反演方差失控，本文的测距窗设置使峰区自然满足可信条件。本文后续的对齐、匹配与融合均以校正后的率域直方图 $F_p(b)\triangleq\hat r_p(b)$ 为输入。

\subsection{单目相对深度引导的局部对齐}
\label{sec:alignment}

设 $p$、$q$ 的单目深度坐标为 $b_p^M$、$b_q^M$，定义单目相对位移
\begin{equation}
  \Delta_{pq}\triangleq b_p^M-b_q^M=(\mu_p-\mu_q)+(e_p-e_q),
  \label{eq:delta}
\end{equation}
其中 $e_p$、$e_q$ 为单目预测误差。$\Delta_{pq}$ 承担两个作用：提供邻居直方图局部搜索中心，并作为 $p$、$q$ 是否属于同一表面的几何证据；其中只含误差之差 $e_p-e_q$，单目全局尺度误差在对齐这一步自然抵消。量化到整数 bin 得粗对齐量 $k_{pq}^{M}=\operatorname{round}(\Delta_{pq})$，以其为中心在窗口 $\mathcal{W}=\{-K,\dots,K\}$ 内对候选偏移 $\delta$ 计算归一化匹配滤波得分
\begin{equation}
  \rho_{pq}(\delta)
  =\NMF\!\big(H_p,\;H_q(b-k_{pq}^{M}-\delta)\big),
  \label{eq:nmf}
\end{equation}
其中 $\NMF(\cdot,\cdot)$ 为归一化匹配滤波算子。取 $\hat\delta_{pq}=\argmax_{\delta\in\mathcal{W}}\rho_{pq}(\delta)$，则实际用于对齐的平移量为
\begin{equation}
  \hat k_{pq}=k_{pq}^{M}+\hat\delta_{pq},
  \qquad
  \widetilde H_{pq}(b)=H_q(b-\hat k_{pq}),
  \label{eq:final_shift}
\end{equation}
其中 $\widetilde H_{pq}$ 为对齐到目标坐标系的邻居直方图；最优位置处的归一化得分记为 $\rho_{pq}=\max_\delta\rho_{pq}(\delta)\in(0,1)$。单目先验只负责将搜索限制在合理局部范围内，最终对齐位置由观测数据的匹配峰决定，$\rho_{pq}$ 记录该峰的匹配质量。

\subsection{可融合性后验与贝叶斯融合权重}
\label{sec:fusibility}

定义二值潜变量 $z_{pq}\in\{0,1\}$：$z_{pq}=1$ 表示邻居 $q$ 可融合（对齐后可用于估计 $\mu_p$），$z_{pq}=0$ 表示不可融合（如来自不同表面或噪声巧合峰）。判断证据为 $(\Delta_{pq},\rho_{pq})$。设可融合状态下 $\Delta_{pq}\mid z_{pq}=1\sim\mathcal{N}(0,\sigma_s^2)$，不可融合状态下 $\Delta_{pq}\mid z_{pq}=0\sim\mathcal{N}(0,\sigma_d^2)$ 且 $\sigma_d>\sigma_s$，则几何似然比为
\begin{equation}
  \Lambda_{\Delta}(\Delta_{pq})
  =\frac{\sigma_d}{\sigma_s}
  \exp\!\left[-\frac{\Delta_{pq}^2}{2}\Big(\frac{1}{\sigma_s^2}-\frac{1}{\sigma_d^2}\Big)\right],
  \label{eq:lambda_delta}
\end{equation}
其中 $\sigma_s$、$\sigma_d$ 分别刻画同一表面内单目局部相对误差的集中程度与跨表面时含深度跳变的散布程度。$\Lambda_\Delta$ 随 $|\Delta_{pq}|$ 单调递减，构成连续的几何软门控，而非硬阈值删除。匹配证据经条件 Beta 分布参数化为似然比 $\Lambda_\rho(\rho_{pq}\mid\Delta_{pq})$（有标注数据时可直接拟合）。由 Bayes 定理，后验可融合概率为
\begin{equation}
  \gamma_{pq}
  =\frac{\pi_{pq}\Lambda_{\Delta}\Lambda_{\rho}}
  {(1-\pi_{pq})+\pi_{pq}\Lambda_{\Delta}\Lambda_{\rho}},
  \label{eq:gamma}
\end{equation}
其中 $\pi_{pq}=\mathbb{P}(z_{pq}=1)$ 为先验可融合概率。等价地在 logit 域中两类证据可加：几何提供软门控，匹配验证实际观测支持。

定义最终对齐误差 $\epsilon_{pq}\triangleq\hat k_{pq}-(\mu_p-\mu_q)$，结合式~\eqref{eq:delta}与式~\eqref{eq:final_shift}得三来源分解
\begin{equation}
  \epsilon_{pq}=(e_p-e_q)+r_{pq}+\hat\delta_{pq},
  \label{eq:epsilon}
\end{equation}
其中 $r_{pq}=k_{pq}^{M}-\Delta_{pq}\in[-\tfrac12,\tfrac12]$ 为取整残差。前两项由先验与离散化决定，第三项由观测选出且与前两者一般不独立，故各方差不能直接相加，对齐方差 $\sigma_{\epsilon,pq}^2$ 由配对标注数据或重复观测估计。对齐直方图的回波中心观测为 $y_{pq}=\mu_p+\epsilon_{pq}+n_{pq}$，其中测量噪声 $n_{pq}$ 的方差 $v_q$ 由式~\eqref{eq:fisher}给出。当测量与对齐误差条件独立且对齐偏差已校准时，总有效方差为
\begin{equation}
  v_{pq}^{\mathrm{eff}}=v_q+\sigma_{\epsilon,pq}^{2}.
  \label{eq:veff}
\end{equation}
将可融合状态建模为高斯观测 $y_{pq}\mid z_{pq}=1\sim\mathcal{N}(\mu_p,v_{pq}^{\mathrm{eff}})$，不可融合状态视为与 $\mu_p$ 无关的异常分量；对完整数据负对数似然关于 $z_{pq}$ 取后验期望（利用 $\mathbb{E}[z_{pq}\mid\Delta_{pq},\rho_{pq}]=\gamma_{pq}$），得到邻居的后验期望精度
\begin{equation}
  c_{pq}=\frac{\gamma_{pq}}{v_{pq}^{\mathrm{eff}}}
  =\frac{\gamma_{pq}}{v_q+\sigma_{\epsilon,pq}^{2}},
  \label{eq:c_pq}
\end{equation}
它将``是否值得融合''（$\gamma_{pq}$）与``融合后有多准确''（分母）统一到同一统计量。令归一化系数 $\alpha_{pq}=c_{pq}/\sum_{r}c_{pr}$，融合估计为 $\hat\mu_p=\sum_{q\in\mathcal{N}_p}\alpha_{pq}y_{pq}$。代入式~\eqref{eq:gamma}展开可见，理想权重呈``可靠性（$\gamma_{pq}$）$\times$ 精度（$1/v_{pq}^{\mathrm{eff}}$）''的因子结构，为工程实现给出明确的分解方向。

\subsection{可标定的工程实现}
\label{sec:engineering}

理想权重要求 $\gamma_{pq}$、$v_q$、$\sigma_{\epsilon,pq}^2$ 在运行时逐像素获得，而单次观测中均不可直接观测；可直接计算的只有 $\Delta_{pq}$、$\hat k_{pq}$、$\rho_{pq}$。本文不为其虚构运行时估计，而是构造与理想似然比同形状的代理，全部自由参数由验证数据离线标定。可靠性因子取
\begin{equation}
  g_{\Delta}(\Delta_{pq})=\exp\!\Big(-\frac{\Delta_{pq}^2}{2\sigma_w^2}\Big),
  \qquad
  g_{\rho}(\rho_{pq})=\rho_{\mathrm{floor}}+(1-\rho_{\mathrm{floor}})\rho_{pq},
  \label{eq:reliability}
\end{equation}
其中 $\sigma_w$ 为几何软门控尺度，标定量级对齐单目局部相对误差（对应 $\sigma_s$）；$\rho_{\mathrm{floor}}\in[0,1)$ 为低分邻居的保底贡献，避免匹配得分偶然偏低时硬性删除仍有先验支持的邻居。精度因子保留显式代理
\begin{equation}
  \hat u_q=\frac{1}{\hat N_q+\hat N_0},
  \qquad
  \hat N_q=\sum_{b\in\Omega_q}\big(H_q(b)-\hat\beta_q\big),
  \label{eq:precision_proxy}
\end{equation}
其中 $\hat N_q$ 为直方图峰显著度，即峰值信号窗口 $\Omega_q$ 内扣除背景估计 $\hat\beta_q$ 后的信号光子数代理，由式~\eqref{eq:v_classical}的反比关系对应理想精度因子 $1/v_q$；$\hat N_0>0$ 为正则化项，对应背景光方差下限，防止低信号邻居权重发散。三个标定参数相互耦合较弱，固定搜索窗口半宽 $K$ 与匹配算子后，在验证集上以深度误差（MAE）为目标逐一维搜索确定并冻结。最终融合权重与归一化系数为
\begin{equation}
  w_{pq}=g_{\Delta}(\Delta_{pq})\,g_{\rho}(\rho_{pq})\,\hat u_q,
  \qquad
  \alpha_{pq}=\frac{w_{pq}}{\sum_{r\in\mathcal{N}_p}w_{pr}},
  \label{eq:weight}
\end{equation}
对对齐后的邻域直方图加权融合
\begin{equation}
  H_p^{\mathrm{fuse}}(b)
  =\sum_{q\in\mathcal{N}_p}\alpha_{pq}\,\widetilde H_{pq}(b),
  \label{eq:fuse}
\end{equation}
再由融合直方图提取回波中心，得到最终深度估计 $\hat\mu_p$。当邻域内信号强度变异有限时，可将 $\hat u_q$ 置为常数，退化为仅可靠性加权。整个流程仅使用单目先验与当前 SPAD 观测即可运行，同时保留了理想贝叶斯模型的核心结构：几何一致性抑制跨表面观测，匹配一致性评价对齐可信度，信号强度压低低光子数邻居的贡献。
 引入。
% 依赖 main.tex 导言区：\DeclareMathOperator*{\argmax}{arg\,max} 与 \DeclareMathOperator{\NMF}{NMF}

\section{方法}
\label{sec:method}

% ================= II.A 观测模型 =================
% 内容取自 method_full.tex 小节〈目标回波与背景光的泊松模型〉与〈背景条件下的回波中心方差〉；标题沿用原 2_1。

\subsection{含背景光的直方图观测模型}
\label{sec:spad}

在光子计数系统中，目标回波光子和背景光子可以近似视为独立泊松到达过程。由独立泊松过程的叠加性质，像素 $p$ 在 bin $b$ 上的总计数满足
\begin{equation}
    H_p(b)
    \sim
    \operatorname{Poisson}\!\left(\lambda_p(b;\mu_p)\right),
    \label{eq:poisson_obs}
\end{equation}
其中
\begin{equation}
    \lambda_p(b;\mu_p)
    =
    N_p h(b-\mu_p)+\beta_p.
    \label{eq:lambda_background}
\end{equation}
两式中，$b$ 为直方图的离散 bin 坐标，$N_p$ 为期望信号光子数，$h(\cdot)$ 为归一化系统脉冲响应（impulse response function, IRF），$\mu_p$ 为真实回波中心坐标，$N_p h(b-\mu_p)$ 描述目标回波的期望光子分布；$\beta_p$ 描述每个 bin 的平均背景光子与暗计数。本文假设当前观测窗口内背景近似平坦；若背景随时间变化，只需将 $\beta_p$ 替换为非负函数 $\beta_p(b)$，以下推导仍成立。对每个 bin，由泊松分布的性质直接有
\begin{equation}
    \E[H_p(b)] = \lambda_p(b;\mu_p),
    \quad
    \Var[H_p(b)] = \lambda_p(b;\mu_p).
    \label{eq:poisson_moments_bg}
\end{equation}
因此，背景光一方面提高直方图基线，另一方面增加每个 bin 的泊松方差，从而降低回波中心估计的统计精度。将回波中心 $\mu_p$ 视为待估参数，含背景泊松模型对应的 Fisher 信息为
\begin{equation}
    \mathcal I_p(\mu_p)
    =
    N_p^2
    \sum_b
    \frac{[h'(b-\mu_p)]^2}
    {N_p h(b-\mu_p)+\beta_p}.
    \label{eq:fisher}
\end{equation}
其中 $h'$ 为 IRF 的导数。由此得回波中心估计的 Cram\'er--Rao 下界为
\begin{equation}
    \Var(\hat\mu_p)
    \geq
    \frac{1}{\mathcal I_p(\mu_p)},
    \label{eq:crlb_bg}
\end{equation}
其中 $\hat\mu_p$ 为由直方图构造的回波中心估计量。在局部渐近条件下，以 Fisher 信息倒数近似回波中心估计方差，记为
\begin{equation}
    v_p
    \approx
    \mathcal I_p(\mu_p)^{-1}.
    \label{eq:v_from_fisher}
\end{equation}
因此，背景光越强，Fisher 信息越低，$v_p$ 越大。特别地，当 $\beta_p\rightarrow0$ 且采用高斯 IRF 时，式~\eqref{eq:fisher}退化为经典的无背景近似
\begin{equation}
    v_p\approx\frac{\sigma_h^2}{N_p},
    \label{eq:v_classical}
\end{equation}
其中 $\sigma_h^2$ 是高斯 IRF 的方差。估计方差与信号光子数成反比这一关系将在第~\ref{sec:engineering}节中用于构造可在线计算的精度代理。


\subsection{死时间畸变与 Coates 校正}
\label{sec:coates}

实际 SPAD 在一个测距周期内至多记录最早到达的一个光子，死时间使原始直方图偏离式~\eqref{eq:poisson_obs}的独立泊松模型。设逐 bin 光子到达率为 $r_p(b)$，累积率 $R_p(b)=\sum_{b'\le b}r_p(b')$，帧内有效周期数为 $M$，则原始直方图在 $b$ 处的累积计数满足精确关系 $\mathrm{cum}(b)=M\big(1-e^{-R_p(b)}\big)$，其解析逆（Coates 校正）为
\begin{equation}
  \hat r_p(b)=\ln\frac{M-\mathrm{cum}(b-1)}{M-\mathrm{cum}(b)},
  \label{eq:coates}
\end{equation}
其中 $\mathrm{cum}(b)=\sum_{b'\le b}H_p^{\mathrm{raw}}(b')$ 为原始直方图累积计数。该反演不含待调参数、对任意强度精确可逆：背景平坦底座反演后恢复常数率，窄峰的形状与中心位置近似保持，故校正不参与对齐量的定义，仅改变直方图取值标度与噪声结构。频率代换引入的统计误差为相对量级 $O(1/\sqrt{M})$ 的随机误差，可并入散粒噪声预算；当幸存周期数过少时反演方差失控，本文的测距窗设置使峰区自然满足可信条件。本文后续的对齐、匹配与融合均以校正后的率域直方图 $F_p(b)\triangleq\hat r_p(b)$ 为输入。

\subsection{单目相对深度引导的局部对齐}
\label{sec:alignment}

设 $p$、$q$ 的单目深度坐标为 $b_p^M$、$b_q^M$，定义单目相对位移
\begin{equation}
  \Delta_{pq}\triangleq b_p^M-b_q^M=(\mu_p-\mu_q)+(e_p-e_q),
  \label{eq:delta}
\end{equation}
其中 $e_p$、$e_q$ 为单目预测误差。$\Delta_{pq}$ 承担两个作用：提供邻居直方图局部搜索中心，并作为 $p$、$q$ 是否属于同一表面的几何证据；其中只含误差之差 $e_p-e_q$，单目全局尺度误差在对齐这一步自然抵消。量化到整数 bin 得粗对齐量 $k_{pq}^{M}=\operatorname{round}(\Delta_{pq})$，以其为中心在窗口 $\mathcal{W}=\{-K,\dots,K\}$ 内对候选偏移 $\delta$ 计算归一化匹配滤波得分
\begin{equation}
  \rho_{pq}(\delta)
  =\NMF\!\big(H_p,\;H_q(b-k_{pq}^{M}-\delta)\big),
  \label{eq:nmf}
\end{equation}
其中 $\NMF(\cdot,\cdot)$ 为归一化匹配滤波算子。取 $\hat\delta_{pq}=\argmax_{\delta\in\mathcal{W}}\rho_{pq}(\delta)$，则实际用于对齐的平移量为
\begin{equation}
  \hat k_{pq}=k_{pq}^{M}+\hat\delta_{pq},
  \qquad
  \widetilde H_{pq}(b)=H_q(b-\hat k_{pq}),
  \label{eq:final_shift}
\end{equation}
其中 $\widetilde H_{pq}$ 为对齐到目标坐标系的邻居直方图；最优位置处的归一化得分记为 $\rho_{pq}=\max_\delta\rho_{pq}(\delta)\in(0,1)$。单目先验只负责将搜索限制在合理局部范围内，最终对齐位置由观测数据的匹配峰决定，$\rho_{pq}$ 记录该峰的匹配质量。


\subsection{单目相对深度引导的局部对齐}
\label{sec:alignment}

设 $p$、$q$ 的单目深度坐标为 $b_p^M$、$b_q^M$，定义单目相对位移
\begin{equation}
  \Delta_{pq}\triangleq b_p^M-b_q^M=(\mu_p-\mu_q)+(e_p-e_q),
  \label{eq:delta}
\end{equation}
其中 $e_p$、$e_q$ 为单目预测误差。$\Delta_{pq}$ 承担两个作用：提供邻居直方图局部搜索中心，并作为 $p$、$q$ 是否属于同一表面的几何证据；其中只含误差之差 $e_p-e_q$，单目全局尺度误差在对齐这一步自然抵消。量化到整数 bin 得粗对齐量 $k_{pq}^{M}=\operatorname{round}(\Delta_{pq})$，以其为中心在窗口 $\mathcal{W}=\{-K,\dots,K\}$ 内对候选偏移 $\delta$ 计算归一化匹配滤波得分
\begin{equation}
  \rho_{pq}(\delta)
  =\NMF\!\big(H_p,\;H_q(b-k_{pq}^{M}-\delta)\big),
  \label{eq:nmf}
\end{equation}
其中 $\NMF(\cdot,\cdot)$ 为归一化匹配滤波算子。取 $\hat\delta_{pq}=\argmax_{\delta\in\mathcal{W}}\rho_{pq}(\delta)$，则实际用于对齐的平移量为
\begin{equation}
  \hat k_{pq}=k_{pq}^{M}+\hat\delta_{pq},
  \qquad
  \widetilde H_{pq}(b)=H_q(b-\hat k_{pq}),
  \label{eq:final_shift}
\end{equation}
其中 $\widetilde H_{pq}$ 为对齐到目标坐标系的邻居直方图；最优位置处的归一化得分记为 $\rho_{pq}=\max_\delta\rho_{pq}(\delta)\in(0,1)$。单目先验只负责将搜索限制在合理局部范围内，最终对齐位置由观测数据的匹配峰决定，$\rho_{pq}$ 记录该峰的匹配质量。
 

\subsection{融合误差与最优融合权重}
\label{sec:error_model}

本节将每个邻居视为独立的观测来源，引入测量与对齐误差，刻画各邻居``应该贡献多少''。由式~\eqref{eq:final_shift}，定义最终对齐误差
\begin{equation}
  \epsilon_{pq}\triangleq\hat k_{pq}-(\mu_p-\mu_q)
  =(e_p-e_q)+r_{pq}+\hat\delta_{pq},
  \label{eq:epsilon}
\end{equation}
其中 $r_{pq}=k_{pq}^{M}-\Delta_{pq}\in[-\tfrac12,\tfrac12]$ 为取整残差。对齐后的邻居回波中心位于 $\mu_q+\hat k_{pq}=\mu_p+\epsilon_{pq}$，其观测为
\begin{equation}
  y_{pq}=\mu_p+\epsilon_{pq}+n_{pq},
  \label{eq:neighbor_obs}
\end{equation}
其中 $n_{pq}$ 为Coates校正后的直方图上回波中心的估计误差，方差由式~\eqref{eq:coates_var}经误差传播给出，记为 $v_q$。对齐误差的方差记为 $\sigma_{\epsilon,pq}^2$。当二者条件独立且对齐偏差已校准时，单个邻居的总不确定度为
\begin{equation}
  v_{pq}^{\mathrm{eff}}=v_q+\sigma_{\epsilon,pq}^{2}.
  \label{eq:veff}
\end{equation}
若各邻居的总不确定度逐对已知，则最优线性融合为逆方差加权，邻居 $q$ 的权重为
\begin{equation}
  c_{pq}=\frac{1}{v_{pq}^{\mathrm{eff}}}
  =\frac{1}{v_q+\sigma_{\epsilon,pq}^{2}}.
  \label{eq:c_pq}
\end{equation}
然而这一前提无法满足，$\sigma_{\epsilon,pq}^{2}$ 的逐对取值依赖未知的真实位移与匹配细节，无法从单次观测直接计算；即便由配对数据整体估计，得到的也只是平均散布，而非``这一对''邻居的对齐可靠性。平均散布至多可充当权重中的基底项，逐对的偏离仍需另行处理。


\subsection{可标定的工程实现}
\label{sec:engineering}

理想权重要求 $\gamma_{pq}$、$v_q$、$\sigma_{\epsilon,pq}^2$ 在运行时逐像素获得，而单次观测中均不可直接观测；可直接计算的只有 $\Delta_{pq}$、$\hat k_{pq}$、$\rho_{pq}$。本文不为其虚构运行时估计，而是构造与理想似然比同形状的代理，全部自由参数由验证数据离线标定。可靠性因子取
\begin{equation}
  g_{\Delta}(\Delta_{pq})=\exp\!\Big(-\frac{\Delta_{pq}^2}{2\sigma_w^2}\Big),
  \qquad
  g_{\rho}(\rho_{pq})=\rho_{\mathrm{floor}}+(1-\rho_{\mathrm{floor}})\rho_{pq},
  \label{eq:reliability}
\end{equation}
其中 $\sigma_w$ 为几何软门控尺度，标定量级对齐单目局部相对误差（对应 $\sigma_s$）；$\rho_{\mathrm{floor}}\in[0,1)$ 为低分邻居的保底贡献，避免匹配得分偶然偏低时硬性删除仍有先验支持的邻居。精度因子保留显式代理
\begin{equation}
  \hat u_q=\frac{1}{\hat N_q+\hat N_0},
  \qquad
  \hat N_q=\sum_{b\in\Omega_q}\big(H_q(b)-\hat\beta_q\big),
  \label{eq:precision_proxy}
\end{equation}
其中 $\hat N_q$ 为直方图峰显著度，即峰值信号窗口 $\Omega_q$ 内扣除背景估计 $\hat\beta_q$ 后的信号光子数代理，由式~\eqref{eq:v_classical}的反比关系对应理想精度因子 $1/v_q$；$\hat N_0>0$ 为正则化项，对应背景光方差下限，防止低信号邻居权重发散。三个标定参数相互耦合较弱，固定搜索窗口半宽 $K$ 与匹配算子后，在验证集上以深度误差（MAE）为目标逐一维搜索确定并冻结。最终融合权重与归一化系数为
\begin{equation}
  w_{pq}=g_{\Delta}(\Delta_{pq})\,g_{\rho}(\rho_{pq})\,\hat u_q,
  \qquad
  \alpha_{pq}=\frac{w_{pq}}{\sum_{r\in\mathcal{G}_p}w_{pr}},
  \label{eq:weight}
\end{equation}
对对齐后的邻域直方图加权融合
\begin{equation}
  H_p^{\mathrm{fuse}}(b)
  =\sum_{q\in\mathcal{G}_p}\alpha_{pq}\,\widetilde H_{pq}(b),
  \label{eq:fuse}
\end{equation}
再由融合直方图提取回波中心，得到最终深度估计 $\hat\mu_p$。当邻域内信号强度变异有限时，可将 $\hat u_q$ 置为常数，退化为仅可靠性加权。整个流程仅使用单目先验与当前 SPAD 观测即可运行，同时保留了理想贝叶斯模型的核心结构：几何一致性抑制跨表面观测，匹配一致性评价对齐可信度，信号强度压低低光子数邻居的贡献。
   % 方法章节（独立文件，便于单独修改）

% ================= 实验章节 =================
% 由 main.tex 抽出，通过 % ================= 实验章节 =================
% 由 main.tex 抽出，通过 % ================= 实验章节 =================
% 由 main.tex 抽出，通过 \input{experiment} 引入。
% 依赖 main.tex 导言区的 amsmath、booktabs、multirow 宏包及方法章节的公式标签（eq:weight、eq:final_shift 等）。

\section{实验}

\subsection{实验配置}
\label{sec:setup}

\subsubsection{仿真数据产生}
仿真数据基于 NYU Depth V2~\cite{nyudepth}的 464 个室内场景生成。激光脉冲为 FWHM $0.6\,\mathrm{ns}$ 的高斯波形，单光子雪崩二极管（SPAD）探测器经 1024-bin 时间数字转换（TDC）量化，bin 宽 $66.9\,\mathrm{ps}$，对应 $[0,10.24]\,\mathrm{m}$ 测距量程；每帧累积 $M=1125$ 个测距周期。像素记录遵循死时间约束（每周期至多记录最早到达的一个光子），背景光与暗计数按泊松过程叠加，与式~\eqref{eq:poisson_obs}的观测模型一致。原生 TDC 分辨率取 $230\times310$；$46\times62$ 分辨率由其 $5\times5$ 邻域聚合得到，单像素信号通量相应提升 25 倍。两分辨率下相同的深度真值图（下采样至对应网格）作为评价参考。

三个工作点仅改变背景照度（$5\times10^3$、$10^4$、$5\times10^4\,\mathrm{lux}$，下文简称 5k/10k/50k），激光信号水平保持不变：以 10k 档为例，230×310 分辨率下单像素平均信号光子数约 3.5，背景与暗计数合计约 68.7。按检测信噪比定义 $\mathrm{SNR}=\mu_s/\sqrt{\mu_s+\mu_b+\mathrm{DCR}}$（$20\log_{10}$ 换算为 dB，对 recorded photons 计入死时间），464 帧全像幅聚合的逐像素 SNR 为 $-5.9\pm8.7$、$-8.6\pm8.9$、$-15.4\pm9.4\,\mathrm{dB}$（46×62 结果一致）。SNR 随背景照度上升而单调下降（背景主导区 $\mathrm{SNR}\propto1/\sqrt{\mu_b}$），三个工作点覆盖了从边缘可用到深度饥饿的运行区间。

\subsubsection{评价指标与实现细节}
全部实验在 464 个场景的完整集合上进行，评价指标为绝对平均误差（MAE）、中位数误差与 inlier 比例（相对误差 $\lvert\hat\mu_p-\mu_p\rvert/\mu_p$ 小于 $0.2\%$、$0.5\%$、$1\%$ 的像素占比），与 CASPI~\cite{caspi}原论文口径一致。单目相对深度由 Depth Anything V2（ViT-S）给出，推理一次后缓存复用。实验平台为 RTX 2080 Ti（PyTorch 与 CuPy）。按 \ref{sec:engineering} 节的标定流程，几何软门控尺度 $\sigma_w=16$ bins，匹配得分下限 $\rho_{\mathrm{floor}}=0.3$，正则化项 $\hat N_0=50$；匹配窗口取 $\pm6$ bins，窗搜索宽度 $2K+1=5$ bins（$K=2$）。上述参数在独立验证集上标定后于全部实验中冻结。

\subsection{消融实验}
\label{sec:ablation}

为量化式~\eqref{eq:weight}中各因子的贡献，构造阶梯式消融序列：L0 为单像素基线（Coates 校正后匹配滤波取峰，无融合）；L1 在 L0 基础上对邻域做等权平均（无对齐）；L2 将等权替换为信号强度权重；L3 引入全轴匹配得分对齐但不用先验（搜索中心不依赖 $\Delta_{pq}$）；L4 先验整数对齐加几何软门控（仅 $g_\Delta$）；L5 加入匹配门与信号强度（$g_\rho\hat u_q$，即式~\eqref{eq:weight}的完整权重但不含窗搜索）；L6 在 L5 基础上加入窗搜索重对齐（式~\eqref{eq:final_shift}，$K=2$）。各级均输出完整深度图并按同一口径评价，结果见表~\ref{tab:ablation230}与表~\ref{tab:ablation46}。

\begin{table}[!t]
\centering
\caption{230×310 分辨率下的消融实验结果（464 个场景全集；MAE 单位为 m，inl@1\% 为相对误差小于 1\% 的像素占比）}
\label{tab:ablation230}\scriptsize\setlength{\tabcolsep}{2.5pt}
\begin{tabular}{clcccccc}
\toprule
\multirow{2}{*}{级别} & \multirow{2}{*}{配置} & \multicolumn{2}{c}{$-5.9\pm8.7$} & \multicolumn{2}{c}{$-8.6\pm8.9$} & \multicolumn{2}{c}{$-15.4\pm9.4$} \\
\cmidrule(lr){3-4}\cmidrule(lr){5-6}\cmidrule(lr){7-8}
 & & MAE & inl@1\% & MAE & inl@1\% & MAE & inl@1\% \\
\midrule
L0 & 单像素（Coates+MF） & 0.716 & 39.8 & 0.860 & 33.5 & 1.727 & 13.2 \\
L1 & 不对齐 + 等权平均 & 0.084 & 62.6 & 0.099 & 62.0 & 0.199 & 57.7 \\
L2 & 不对齐 + 信号强度权重 & 0.082 & 62.7 & 0.098 & 62.1 & 0.196 & 58.0 \\
L3 & 无先验 + 全轴匹配得分对齐 & 1.571 & 29.5 & 1.653 & 25.4 & 2.055 & 12.2 \\
L4 & 先验整数对齐 + 几何门 & 0.064 & 64.3 & 0.080 & 63.3 & 0.194 & 57.1 \\
L5 & L4 + 匹配门 + 信号强度 & 0.061 & 65.1 & 0.077 & 64.0 & 0.190 & 57.7 \\
L6 & L5 + 窗搜索重对齐 & \textbf{0.058} & \textbf{64.3} & \textbf{0.074} & \textbf{63.9} & \textbf{0.186} & \textbf{57.6} \\
\bottomrule
\end{tabular}
\end{table}

\begin{table}[!t]
\centering
\caption{46×62 分辨率下的消融实验结果（464 个场景全集；MAE 单位为 m，inl@1\% 为相对误差小于 1\% 的像素占比）}
\label{tab:ablation46}\scriptsize\setlength{\tabcolsep}{2.5pt}
\begin{tabular}{clcccccc}
\toprule
\multirow{2}{*}{级别} & \multirow{2}{*}{配置} & \multicolumn{2}{c}{$-5.7\pm8.4$} & \multicolumn{2}{c}{$-8.4\pm8.6$} & \multicolumn{2}{c}{$-15.2\pm9.0$} \\
\cmidrule(lr){3-4}\cmidrule(lr){5-6}\cmidrule(lr){7-8}
 & & MAE & inl@1\% & MAE & inl@1\% & MAE & inl@1\% \\
\midrule
L0 & 单像素（Coates+MF） & 0.665 & 37.0 & 0.814 & 30.9 & 1.695 & 12.5 \\
L1 & 不对齐 + 等权平均 & 0.323 & 22.7 & 0.349 & 22.3 & 0.522 & 19.7 \\
L2 & 不对齐 + 信号强度权重 & 0.317 & 23.0 & 0.342 & 22.5 & 0.519 & 20.0 \\
L3 & 无先验 + 全轴匹配得分对齐 & 1.470 & 27.3 & 1.533 & 23.2 & 2.031 & 10.0 \\
L4 & 先验整数对齐 + 几何门 & 0.214 & 29.3 & 0.235 & 30.1 & 0.460 & 24.1 \\
L5 & L4 + 匹配门 + 信号强度 & 0.197 & 32.4 & 0.221 & 32.5 & 0.451 & 25.2 \\
L6 & L5 + 窗搜索重对齐 & \textbf{0.184} & \textbf{40.1} & \textbf{0.207} & \textbf{38.2} & \textbf{0.427} & \textbf{27.5} \\
\bottomrule
\end{tabular}
\end{table}

结果呈现三个一致的规律。其一，无先验的数据驱动对齐（L3）在两个分辨率下均系统性崩溃（MAE $1.5\sim2.1\,\mathrm{m}$，甚至劣于不融合的 L0）：归一化匹配得分在背景主导的直方图上频繁锁定背景巧合峰，且跨邻居的错位独立发生、无法通过平均抵消。将搜索中心约束为先验整数对齐（L4）后误差立即恢复一个数量级以上，验证了 \ref{sec:alignment} 节“单目先验提供搜索范围、数据决定最终对齐”的分工必要性；对照式~\eqref{eq:delta}，此时起作用的仅是单目相对位移的几何趋势，其全局尺度误差不参与。其二，几何门（L4）与匹配门、信号强度（L5）的逐级引入在两个分辨率下均带来单调改善，说明可靠性因子 $g_\Delta$、$g_\rho$ 与精度因子 $\hat u_q$ 的贡献可分离且方向一致。其三，窗搜索（L6）在统一 $K=2$ 配置下于两个分辨率、三个照度档全部取得最优 MAE 与中位数误差，相对 L5 的净增益为 $3\%\sim7\%$；同时 inl@1\% 不降，表明小窗重对齐在修正邻居残余错位的同时未引入背景巧合峰。

值得指出的是 MAE 与 inl@1\% 在 L0$\to$L1 处的背离：46×62 上等权平均使 MAE 从 $0.665$ 降至 $0.323$，inl@1\% 却从 $37.0\%$ 跌至 $22.7\%$。这是误差分布重组的结果——单像素估计呈双峰分布（约 37\% 的像素亚厘米级准确，其余误差约 $1\,\mathrm{m}$ 量级），等权平均将高信噪比像素与跨边缘、饥饿邻居无差别混合，以头部精度换取尾部均值；后续几何门与匹配门逐级将该头部精度找回，至 L6 时 inl@1\% 反超 L0（$40.1\%$ vs $37.0\%$），实现两端兼顾。这一现象也说明单靠 MAE 会高估朴素融合的价值。

\subsection{对比实验}
\label{sec:comparison}

将本文方法与四类基线在完全相同的 464 个场景、同一真值与同一评价脚本下比较：（1）Argmax，原始直方图直接取峰（无 Coates 校正与匹配滤波）；（2）Raw+MF，原始直方图经匹配滤波后取峰；（3）Coates+MF，即消融实验中的 L0；（4）CASPI~\cite{caspi}，使用官方实现与原论文参数；（5）本文方法（L6，$K=2$）。结果见表~\ref{tab:comp230}与表~\ref{tab:comp46}。

\begin{table}[!t]
\centering
\caption{230×310 分辨率下与基线方法的定量对比（464 个场景全集；上半部分为 MAE，单位 m；下半部分为 inl@1\%，单位 \%；列组标题为全像幅平均检测 SNR，mean$\pm\mathrm{std}$，单位 dB）}
\label{tab:comp230}
\scriptsize\setlength{\tabcolsep}{2.5pt}
\begin{tabular}{lcccccc}
\toprule
\multirow{2}{*}{方法} & \multicolumn{2}{c}{$-5.9\pm8.7$} & \multicolumn{2}{c}{$-8.6\pm8.9$} & \multicolumn{2}{c}{$-15.4\pm9.4$} \\
\cmidrule(lr){2-3}\cmidrule(lr){4-5}\cmidrule(lr){6-7}
 & MAE & inl@1\% & MAE & inl@1\% & MAE & inl@1\% \\
\midrule
\multicolumn{7}{l}{MAE (m)$\;\downarrow$}\\
Argmax & 1.645 & 1.728 & 1.980 \\
Raw+MF & 0.415 & 0.559 & 1.073 \\
Coates+MF & 0.716 & 0.860 & 1.727 \\
CASPI~\cite{caspi} & 0.097 & 0.140 & 0.375 \\
Ours & \textbf{0.058} & \textbf{0.074} & \textbf{0.186} \\
\midrule
\multicolumn{7}{l}{inl@1\% $\uparrow$}\\
Argmax & 18.3 & 14.8 & 7.2 \\
Raw+MF & 43.1 & 36.3 & 20.2 \\
Coates+MF & 39.8 & 33.5 & 13.2 \\
CASPI~\cite{caspi} & 53.4 & 51.7 & 43.6 \\
Ours & \textbf{64.3} & \textbf{63.9} & \textbf{57.6} \\
\bottomrule
\end{tabular}
\end{table}

\begin{table}[!t]
\centering
\caption{46×62 分辨率下与基线方法的定量对比（464 个场景全集；上半部分为 MAE，单位 m；下半部分为 inl@1\%，单位 \%；列组标题为全像幅平均检测 SNR，mean$\pm\mathrm{std}$，单位 dB）}
\label{tab:comp46}
\scriptsize\setlength{\tabcolsep}{2.5pt}
\begin{tabular}{lcccccc}
\toprule
\multirow{2}{*}{方法} & \multicolumn{2}{c}{$-5.7\pm8.4$} & \multicolumn{2}{c}{$-8.4\pm8.6$} & \multicolumn{2}{c}{$-15.2\pm9.0$} \\
\cmidrule(lr){2-3}\cmidrule(lr){4-5}\cmidrule(lr){6-7}
 & MAE & inl@1\% & MAE & inl@1\% & MAE & inl@1\% \\
\midrule
\multicolumn{7}{l}{MAE (m)$\;\downarrow$}\\
Argmax & 1.645 & 1.724 & 1.976 \\
Raw+MF & 0.404 & 0.553 & 1.048 \\
Coates+MF & 0.665 & 0.814 & 1.695 \\
CASPI~\cite{caspi} & 0.207 & 0.282 & 0.648 \\
Ours & \textbf{0.184} & \textbf{0.207} & \textbf{0.427} \\
\midrule
\multicolumn{7}{l}{inl@1\% $\uparrow$}\\
Argmax & 18.3 & 14.9 & 7.2 \\
Raw+MF & 39.6 & 33.1 & 18.1 \\
Coates+MF & 37.0 & 30.9 & 12.5 \\
CASPI~\cite{caspi} & 38.1 & 35.0 & 25.4 \\
Ours & \textbf{40.1} & \textbf{38.2} & \textbf{27.5} \\
\bottomrule
\end{tabular}
\end{table}

对比给出三点观察。首先，传统预处理的选择比直觉敏感：Coates 校正作为反卷积类操作，在光子饥饿条件下将稀疏计数中的噪声放大为高频伪峰，使 Coates+MF 系统性劣于不校正的 Raw+MF（MAE 劣化 $47\%\sim73\%$），匹配滤波在此承担的实为抑制该校正副作用的角色；该差异在 argmax 级同样存在（不校正直方图直接取峰优于 Coates 后取峰 2 倍以上），说明这是校正本身的统计性质而非某一实现的问题。其次，CASPI 在 46×62 上是强基线（较单像素基线改善 $55\%\sim69\%$），但在 230×310 上相对自身退化、与本文方法的差距拉大至 $40\%\sim51\%$：其非局部校正依赖像素间统计相似性，原生分辨率下单像素光子数下降一个量级后相似性证据随之减弱，而本文方法因单目先验提供了与光子统计无关的几何约束，受分辨率与照度的影响显著更小。最后，在全部 12 个工作点中本文方法均取得最优 MAE 与 inl@1\%，且尾部指标的领先幅度大于均值指标（230×310 下 inl@1\% 较 CASPI 高 $6\sim14$ 个百分点），与消融实验中融合权重对不可融合邻居的抑制作用一致。效率方面，在相同的 RTX 2080 Ti 上，CASPI 处理单帧需 89.6 s（230×310），本文完整融合流水线（含单目推理）约 55 s，且该耗时同时完成了全部八个消融变体的计算。
 引入。
% 依赖 main.tex 导言区的 amsmath、booktabs、multirow 宏包及方法章节的公式标签（eq:weight、eq:final_shift 等）。

\section{实验}

\subsection{实验配置}
\label{sec:setup}

\subsubsection{仿真数据产生}
仿真数据基于 NYU Depth V2~\cite{nyudepth}的 464 个室内场景生成。激光脉冲为 FWHM $0.6\,\mathrm{ns}$ 的高斯波形，单光子雪崩二极管（SPAD）探测器经 1024-bin 时间数字转换（TDC）量化，bin 宽 $66.9\,\mathrm{ps}$，对应 $[0,10.24]\,\mathrm{m}$ 测距量程；每帧累积 $M=1125$ 个测距周期。像素记录遵循死时间约束（每周期至多记录最早到达的一个光子），背景光与暗计数按泊松过程叠加，与式~\eqref{eq:poisson_obs}的观测模型一致。原生 TDC 分辨率取 $230\times310$；$46\times62$ 分辨率由其 $5\times5$ 邻域聚合得到，单像素信号通量相应提升 25 倍。两分辨率下相同的深度真值图（下采样至对应网格）作为评价参考。

三个工作点仅改变背景照度（$5\times10^3$、$10^4$、$5\times10^4\,\mathrm{lux}$，下文简称 5k/10k/50k），激光信号水平保持不变：以 10k 档为例，230×310 分辨率下单像素平均信号光子数约 3.5，背景与暗计数合计约 68.7。按检测信噪比定义 $\mathrm{SNR}=\mu_s/\sqrt{\mu_s+\mu_b+\mathrm{DCR}}$（$20\log_{10}$ 换算为 dB，对 recorded photons 计入死时间），464 帧全像幅聚合的逐像素 SNR 为 $-5.9\pm8.7$、$-8.6\pm8.9$、$-15.4\pm9.4\,\mathrm{dB}$（46×62 结果一致）。SNR 随背景照度上升而单调下降（背景主导区 $\mathrm{SNR}\propto1/\sqrt{\mu_b}$），三个工作点覆盖了从边缘可用到深度饥饿的运行区间。

\subsubsection{评价指标与实现细节}
全部实验在 464 个场景的完整集合上进行，评价指标为绝对平均误差（MAE）、中位数误差与 inlier 比例（相对误差 $\lvert\hat\mu_p-\mu_p\rvert/\mu_p$ 小于 $0.2\%$、$0.5\%$、$1\%$ 的像素占比），与 CASPI~\cite{caspi}原论文口径一致。单目相对深度由 Depth Anything V2（ViT-S）给出，推理一次后缓存复用。实验平台为 RTX 2080 Ti（PyTorch 与 CuPy）。按 \ref{sec:engineering} 节的标定流程，几何软门控尺度 $\sigma_w=16$ bins，匹配得分下限 $\rho_{\mathrm{floor}}=0.3$，正则化项 $\hat N_0=50$；匹配窗口取 $\pm6$ bins，窗搜索宽度 $2K+1=5$ bins（$K=2$）。上述参数在独立验证集上标定后于全部实验中冻结。

\subsection{消融实验}
\label{sec:ablation}

为量化式~\eqref{eq:weight}中各因子的贡献，构造阶梯式消融序列：L0 为单像素基线（Coates 校正后匹配滤波取峰，无融合）；L1 在 L0 基础上对邻域做等权平均（无对齐）；L2 将等权替换为信号强度权重；L3 引入全轴匹配得分对齐但不用先验（搜索中心不依赖 $\Delta_{pq}$）；L4 先验整数对齐加几何软门控（仅 $g_\Delta$）；L5 加入匹配门与信号强度（$g_\rho\hat u_q$，即式~\eqref{eq:weight}的完整权重但不含窗搜索）；L6 在 L5 基础上加入窗搜索重对齐（式~\eqref{eq:final_shift}，$K=2$）。各级均输出完整深度图并按同一口径评价，结果见表~\ref{tab:ablation230}与表~\ref{tab:ablation46}。

\begin{table}[!t]
\centering
\caption{230×310 分辨率下的消融实验结果（464 个场景全集；MAE 单位为 m，inl@1\% 为相对误差小于 1\% 的像素占比）}
\label{tab:ablation230}\scriptsize\setlength{\tabcolsep}{2.5pt}
\begin{tabular}{clcccccc}
\toprule
\multirow{2}{*}{级别} & \multirow{2}{*}{配置} & \multicolumn{2}{c}{$-5.9\pm8.7$} & \multicolumn{2}{c}{$-8.6\pm8.9$} & \multicolumn{2}{c}{$-15.4\pm9.4$} \\
\cmidrule(lr){3-4}\cmidrule(lr){5-6}\cmidrule(lr){7-8}
 & & MAE & inl@1\% & MAE & inl@1\% & MAE & inl@1\% \\
\midrule
L0 & 单像素（Coates+MF） & 0.716 & 39.8 & 0.860 & 33.5 & 1.727 & 13.2 \\
L1 & 不对齐 + 等权平均 & 0.084 & 62.6 & 0.099 & 62.0 & 0.199 & 57.7 \\
L2 & 不对齐 + 信号强度权重 & 0.082 & 62.7 & 0.098 & 62.1 & 0.196 & 58.0 \\
L3 & 无先验 + 全轴匹配得分对齐 & 1.571 & 29.5 & 1.653 & 25.4 & 2.055 & 12.2 \\
L4 & 先验整数对齐 + 几何门 & 0.064 & 64.3 & 0.080 & 63.3 & 0.194 & 57.1 \\
L5 & L4 + 匹配门 + 信号强度 & 0.061 & 65.1 & 0.077 & 64.0 & 0.190 & 57.7 \\
L6 & L5 + 窗搜索重对齐 & \textbf{0.058} & \textbf{64.3} & \textbf{0.074} & \textbf{63.9} & \textbf{0.186} & \textbf{57.6} \\
\bottomrule
\end{tabular}
\end{table}

\begin{table}[!t]
\centering
\caption{46×62 分辨率下的消融实验结果（464 个场景全集；MAE 单位为 m，inl@1\% 为相对误差小于 1\% 的像素占比）}
\label{tab:ablation46}\scriptsize\setlength{\tabcolsep}{2.5pt}
\begin{tabular}{clcccccc}
\toprule
\multirow{2}{*}{级别} & \multirow{2}{*}{配置} & \multicolumn{2}{c}{$-5.7\pm8.4$} & \multicolumn{2}{c}{$-8.4\pm8.6$} & \multicolumn{2}{c}{$-15.2\pm9.0$} \\
\cmidrule(lr){3-4}\cmidrule(lr){5-6}\cmidrule(lr){7-8}
 & & MAE & inl@1\% & MAE & inl@1\% & MAE & inl@1\% \\
\midrule
L0 & 单像素（Coates+MF） & 0.665 & 37.0 & 0.814 & 30.9 & 1.695 & 12.5 \\
L1 & 不对齐 + 等权平均 & 0.323 & 22.7 & 0.349 & 22.3 & 0.522 & 19.7 \\
L2 & 不对齐 + 信号强度权重 & 0.317 & 23.0 & 0.342 & 22.5 & 0.519 & 20.0 \\
L3 & 无先验 + 全轴匹配得分对齐 & 1.470 & 27.3 & 1.533 & 23.2 & 2.031 & 10.0 \\
L4 & 先验整数对齐 + 几何门 & 0.214 & 29.3 & 0.235 & 30.1 & 0.460 & 24.1 \\
L5 & L4 + 匹配门 + 信号强度 & 0.197 & 32.4 & 0.221 & 32.5 & 0.451 & 25.2 \\
L6 & L5 + 窗搜索重对齐 & \textbf{0.184} & \textbf{40.1} & \textbf{0.207} & \textbf{38.2} & \textbf{0.427} & \textbf{27.5} \\
\bottomrule
\end{tabular}
\end{table}

结果呈现三个一致的规律。其一，无先验的数据驱动对齐（L3）在两个分辨率下均系统性崩溃（MAE $1.5\sim2.1\,\mathrm{m}$，甚至劣于不融合的 L0）：归一化匹配得分在背景主导的直方图上频繁锁定背景巧合峰，且跨邻居的错位独立发生、无法通过平均抵消。将搜索中心约束为先验整数对齐（L4）后误差立即恢复一个数量级以上，验证了 \ref{sec:alignment} 节“单目先验提供搜索范围、数据决定最终对齐”的分工必要性；对照式~\eqref{eq:delta}，此时起作用的仅是单目相对位移的几何趋势，其全局尺度误差不参与。其二，几何门（L4）与匹配门、信号强度（L5）的逐级引入在两个分辨率下均带来单调改善，说明可靠性因子 $g_\Delta$、$g_\rho$ 与精度因子 $\hat u_q$ 的贡献可分离且方向一致。其三，窗搜索（L6）在统一 $K=2$ 配置下于两个分辨率、三个照度档全部取得最优 MAE 与中位数误差，相对 L5 的净增益为 $3\%\sim7\%$；同时 inl@1\% 不降，表明小窗重对齐在修正邻居残余错位的同时未引入背景巧合峰。

值得指出的是 MAE 与 inl@1\% 在 L0$\to$L1 处的背离：46×62 上等权平均使 MAE 从 $0.665$ 降至 $0.323$，inl@1\% 却从 $37.0\%$ 跌至 $22.7\%$。这是误差分布重组的结果——单像素估计呈双峰分布（约 37\% 的像素亚厘米级准确，其余误差约 $1\,\mathrm{m}$ 量级），等权平均将高信噪比像素与跨边缘、饥饿邻居无差别混合，以头部精度换取尾部均值；后续几何门与匹配门逐级将该头部精度找回，至 L6 时 inl@1\% 反超 L0（$40.1\%$ vs $37.0\%$），实现两端兼顾。这一现象也说明单靠 MAE 会高估朴素融合的价值。

\subsection{对比实验}
\label{sec:comparison}

将本文方法与四类基线在完全相同的 464 个场景、同一真值与同一评价脚本下比较：（1）Argmax，原始直方图直接取峰（无 Coates 校正与匹配滤波）；（2）Raw+MF，原始直方图经匹配滤波后取峰；（3）Coates+MF，即消融实验中的 L0；（4）CASPI~\cite{caspi}，使用官方实现与原论文参数；（5）本文方法（L6，$K=2$）。结果见表~\ref{tab:comp230}与表~\ref{tab:comp46}。

\begin{table}[!t]
\centering
\caption{230×310 分辨率下与基线方法的定量对比（464 个场景全集；上半部分为 MAE，单位 m；下半部分为 inl@1\%，单位 \%；列组标题为全像幅平均检测 SNR，mean$\pm\mathrm{std}$，单位 dB）}
\label{tab:comp230}
\scriptsize\setlength{\tabcolsep}{2.5pt}
\begin{tabular}{lcccccc}
\toprule
\multirow{2}{*}{方法} & \multicolumn{2}{c}{$-5.9\pm8.7$} & \multicolumn{2}{c}{$-8.6\pm8.9$} & \multicolumn{2}{c}{$-15.4\pm9.4$} \\
\cmidrule(lr){2-3}\cmidrule(lr){4-5}\cmidrule(lr){6-7}
 & MAE & inl@1\% & MAE & inl@1\% & MAE & inl@1\% \\
\midrule
\multicolumn{7}{l}{MAE (m)$\;\downarrow$}\\
Argmax & 1.645 & 1.728 & 1.980 \\
Raw+MF & 0.415 & 0.559 & 1.073 \\
Coates+MF & 0.716 & 0.860 & 1.727 \\
CASPI~\cite{caspi} & 0.097 & 0.140 & 0.375 \\
Ours & \textbf{0.058} & \textbf{0.074} & \textbf{0.186} \\
\midrule
\multicolumn{7}{l}{inl@1\% $\uparrow$}\\
Argmax & 18.3 & 14.8 & 7.2 \\
Raw+MF & 43.1 & 36.3 & 20.2 \\
Coates+MF & 39.8 & 33.5 & 13.2 \\
CASPI~\cite{caspi} & 53.4 & 51.7 & 43.6 \\
Ours & \textbf{64.3} & \textbf{63.9} & \textbf{57.6} \\
\bottomrule
\end{tabular}
\end{table}

\begin{table}[!t]
\centering
\caption{46×62 分辨率下与基线方法的定量对比（464 个场景全集；上半部分为 MAE，单位 m；下半部分为 inl@1\%，单位 \%；列组标题为全像幅平均检测 SNR，mean$\pm\mathrm{std}$，单位 dB）}
\label{tab:comp46}
\scriptsize\setlength{\tabcolsep}{2.5pt}
\begin{tabular}{lcccccc}
\toprule
\multirow{2}{*}{方法} & \multicolumn{2}{c}{$-5.7\pm8.4$} & \multicolumn{2}{c}{$-8.4\pm8.6$} & \multicolumn{2}{c}{$-15.2\pm9.0$} \\
\cmidrule(lr){2-3}\cmidrule(lr){4-5}\cmidrule(lr){6-7}
 & MAE & inl@1\% & MAE & inl@1\% & MAE & inl@1\% \\
\midrule
\multicolumn{7}{l}{MAE (m)$\;\downarrow$}\\
Argmax & 1.645 & 1.724 & 1.976 \\
Raw+MF & 0.404 & 0.553 & 1.048 \\
Coates+MF & 0.665 & 0.814 & 1.695 \\
CASPI~\cite{caspi} & 0.207 & 0.282 & 0.648 \\
Ours & \textbf{0.184} & \textbf{0.207} & \textbf{0.427} \\
\midrule
\multicolumn{7}{l}{inl@1\% $\uparrow$}\\
Argmax & 18.3 & 14.9 & 7.2 \\
Raw+MF & 39.6 & 33.1 & 18.1 \\
Coates+MF & 37.0 & 30.9 & 12.5 \\
CASPI~\cite{caspi} & 38.1 & 35.0 & 25.4 \\
Ours & \textbf{40.1} & \textbf{38.2} & \textbf{27.5} \\
\bottomrule
\end{tabular}
\end{table}

对比给出三点观察。首先，传统预处理的选择比直觉敏感：Coates 校正作为反卷积类操作，在光子饥饿条件下将稀疏计数中的噪声放大为高频伪峰，使 Coates+MF 系统性劣于不校正的 Raw+MF（MAE 劣化 $47\%\sim73\%$），匹配滤波在此承担的实为抑制该校正副作用的角色；该差异在 argmax 级同样存在（不校正直方图直接取峰优于 Coates 后取峰 2 倍以上），说明这是校正本身的统计性质而非某一实现的问题。其次，CASPI 在 46×62 上是强基线（较单像素基线改善 $55\%\sim69\%$），但在 230×310 上相对自身退化、与本文方法的差距拉大至 $40\%\sim51\%$：其非局部校正依赖像素间统计相似性，原生分辨率下单像素光子数下降一个量级后相似性证据随之减弱，而本文方法因单目先验提供了与光子统计无关的几何约束，受分辨率与照度的影响显著更小。最后，在全部 12 个工作点中本文方法均取得最优 MAE 与 inl@1\%，且尾部指标的领先幅度大于均值指标（230×310 下 inl@1\% 较 CASPI 高 $6\sim14$ 个百分点），与消融实验中融合权重对不可融合邻居的抑制作用一致。效率方面，在相同的 RTX 2080 Ti 上，CASPI 处理单帧需 89.6 s（230×310），本文完整融合流水线（含单目推理）约 55 s，且该耗时同时完成了全部八个消融变体的计算。
 引入。
% 依赖 main.tex 导言区的 amsmath、booktabs、multirow 宏包及方法章节的公式标签（eq:weight、eq:final_shift 等）。

\section{实验}

\subsection{实验配置}
\label{sec:setup}

\subsubsection{仿真数据产生}
仿真采用 Scholes 等人~\cite{Scholes2023}提出的单光子深度成像数值仿真方法：在给定光学系统与采集时间的条件下，对脉冲激光传输衰减、场景反射、背景光注入与 SPAD 探测器的逐光子响应进行建模，直接输出含死时间畸变的原始光子计数直方图。本文完整沿用其系统配置，全部参数列于表~\ref{tab:simcfg}。SPAD 阵列密度取两档，即高分辨率 $230\times310$ 与低分辨率 $46\times62$，两档网格下光学与探测器参数完全相同。背景光照度取三档：$5\times10^{3}$、$1\times10^{4}$、$5\times10^{4}\,\mathrm{lux}$。为量化工作点难度，定义逐像素检测信噪比
\begin{equation}
  \mathrm{SNR}_p=\frac{N^{\mathrm{sig}}_p}{\sqrt{N^{\mathrm{sig}}_p+N^{\mathrm{bkg}}_p}},
  \qquad
  \mathrm{SNR}_p^{\,\mathrm{dB}}=20\log_{10}\mathrm{SNR}_p,
  \label{eq:snr}
\end{equation}
其中 $N^{\mathrm{sig}}_p$、$N^{\mathrm{bkg}}_p$ 分别为像素 $p$ 单帧内期望记录的信号事件数与噪声事件数，由仿真器在死时间竞争下解析分解得到。背景主导（$N^{\mathrm{bkg}}_p\gg N^{\mathrm{sig}}_p$）时，$\mathrm{SNR}$ 随背景照度按 $1/\sqrt{N^{\mathrm{bkg}}_p}$ 下降。

\begin{table}[!t]
\centering
\caption{LiDAR 仿真器系统配置}
\label{tab:simcfg}\small\setlength{\tabcolsep}{15pt}
\begin{tabular}{@{}cll@{}}
\toprule
模块 & 参数 & 数值 \\
\midrule
\multirow{6}{*}{发射端} & 激光中心波长 & $671\,\mathrm{nm}$ \\
 & 激光脉冲半高全宽 & $0.6\,\mathrm{ns}$ \\
 & 单脉冲激光能量 & $1.219\times10^{-10}\,\mathrm{J}$ \\
 & 激光脉冲重复频率 & $2.25\,\mathrm{MHz}$ \\
 & 照明光纤芯径 & $550\,\mu\mathrm{m}$ \\
 & 照明透镜焦距 & $8\,\mathrm{mm}$ \\
\midrule
\multirow{10}{*}{接收端} & 接收镜头光圈数 & $2.0$ \\
 & 光子探测效率 & $18\%$ \\
 & 有效像元尺寸 & $9.2\times9.2\,\mu\mathrm{m}^2$ \\
 & 暗计数率 & $126\,\mathrm{s^{-1}}$ \\
 & 激光脉冲时间抖动 & $220\,\mathrm{ps}$ \\
 & 仿真精细时间步长 & $50\,\mathrm{ps}$ \\
 & 直方图 bin 数 & $1024$ \\
 & 测距量程 & $[0,\,10.24]\,\mathrm{m}$ \\
 & 帧率 & $2000\,\mathrm{fps}$ \\
 & 每帧累积测距周期数 & $1125$ \\
\bottomrule
\end{tabular}
\end{table}

仿真场景取 NYU Depth V2 Labeled Dataset~\cite{Silberman2012}全部 $464$ 个室内场景的第一帧。每帧 $640\times480$ 图像先裁剪四周 $10$ 像素边缘，再经面积插值降采样至目标分辨率；深度真值同步降采样至对应分辨率，作为评价参考。表~\ref{tab:snr}汇总了全部 464 个场景在两种分辨率、三档照度下逐像素 SNR 统计。三个工作点覆盖了从勉强可用到深度饥饿的运行区间：5k lux 档尚有约四分之一像素 SNR 高于 $0\,\mathrm{dB}$；50k lux 档中位数低至 $-15.6\,\mathrm{dB}$，即便最好的 5\% 像素也仅勉强越过 $0\,\mathrm{dB}$。$\mathrm{SNR}<0\,\mathrm{dB}$ 即期望信号计数低于单帧总计数的标准差，信号峰淹没于统计噪声之中，此类像素在 50k lux 档占比超过 $94\%$。两种分辨率的统计高度接近，均值差约 $0.2\,\mathrm{dB}$。

\begin{table}[!t]
\centering
\caption{全部 464 个场景的逐像素检测 SNR 统计（单位 dB）}
\label{tab:snr}\small\setlength{\tabcolsep}{2pt}
\begin{tabular}{ccccccc}
\toprule
分辨率 & 照度 & mean$\pm$std & 中位数 & P5 & P95 & $<0$\,dB 占比 \\
\midrule
\multirow{3}{*}{$230\times310$} & 5k & $-5.9\pm8.7$ & $-6.0$ & $-20.3$ & $8.4$ & 75.4\% \\
 & 10k & $-8.6\pm8.9$ & $-8.7$ & $-23.2$ & $6.3$ & 83.3\% \\
 & 50k & $-15.4\pm9.4$ & $-15.6$ & $-30.6$ & $0.5$ & 94.5\% \\
\midrule
\multirow{3}{*}{$46\times62$} & 5k & $-5.7\pm8.4$ & $-5.9$ & $-19.6$ & $8.4$ & 75.3\% \\
 & 10k & $-8.4\pm8.6$ & $-8.6$ & $-22.5$ & $6.3$ & 83.3\% \\
 & 50k & $-15.2\pm9.0$ & $-15.5$ & $-29.9$ & $0.5$ & 94.5\% \\
\bottomrule
\end{tabular}
\end{table}

\subsubsection{实现细节与评价指标}
本文方法的全部自由参数汇总于表~\ref{tab:impl}，符号与第~\ref{sec:method}节一致。按第~\ref{sec:engineering}节的标定流程，上述参数以 MAE 为目标在验证集上逐维搜索确定：验证集由随机抽取的 64 个场景构成，标定在低分辨率数据上进行；所得参数在其余场景、原生分辨率与全部照度工作点中沿用同一组取值并全程冻结。标定过程中未观察到对参数取值的显著敏感。单目相对深度由 Depth-Anything-V2-Small~\cite{Yang2024}给出。评价指标包括平均绝对误差（MAE）、中位数误差与 inlier 比例，其中 inlier 比例为相对误差 $\lvert\hat d_p-d_p\rvert/d_p$ 小于 $0.2\%$、$0.5\%$、$1\%$ 的像素占比，$\hat d_p$、$d_p$ 分别为估计深度与真值深度。

\begin{table}[!t]
\centering
\caption{本文方法的自由参数取值}
\label{tab:impl}\small\setlength{\tabcolsep}{1pt}
\begin{tabular}{@{}cll@{}}
\toprule
符号 & 含义 & 取值 \\
\midrule
$R$ & 邻域 $\mathcal{G}_p$ 半径（共 $(2R{+}1)^2$ 个邻居） & $6$ \\
$\sigma_{\mathrm{mf}}$ & 取峰匹配滤波核标准差（$\ge\sigma_h$） & $6.73$ bins \\
$K$ & 窗搜索半宽（$2K{+}1$ 个候选偏移） & $5$ \\
$\sigma_w$ & 几何软门控尺度 & $16$ bins \\
$\rho_{\mathrm{floor}}$ & 匹配得分保底贡献 & $0.3$ \\
$\hat N_0$ & 对齐残差强度标定量纲（$\sigma_h^2/\sigma_\epsilon^2$） & $50$ \\
$\Omega_q$ & 峰值信号窗口半宽 & $27$ bins \\
$\hat\beta_q$ & 背景估计 & 直方图 $20\%$ 分位数 \\
\bottomrule
\end{tabular}
\end{table}


\subsection{消融实验}
\label{sec:ablation}

为量化式~\eqref{eq:weight}中各因子的贡献，构造阶梯式消融序列：L0 为单像素基线（Coates 校正后匹配滤波取峰，无融合）；L1 在 L0 基础上对邻域做等权平均（无对齐）；L2 将等权替换为信号强度权重；L3 引入全轴匹配得分对齐但不用先验（搜索中心不依赖 $\Delta_{pq}$）；L4 先验整数对齐加几何软门控（仅 $g_\Delta$）；L5 加入匹配门与幅度响应（即式~\eqref{eq:weight}的完整权重但不含窗搜索）；L6 在 L5 基础上加入窗搜索重对齐（式~\eqref{eq:final_shift}，$K=5$）。各级均输出完整深度图并按同一口径评价，结果见表~\ref{tab:ablation230}与表~\ref{tab:ablation46}。

\begin{table}[!t]
\centering
\caption{230×310 分辨率下的消融实验结果（464 个场景全集；MAE 单位为 m，inl@1\% 为相对误差小于 1\% 的像素占比）}
\label{tab:ablation230}\scriptsize\setlength{\tabcolsep}{2.5pt}
\begin{tabular}{clcccccc}
\toprule
\multirow{2}{*}{级别} & \multirow{2}{*}{配置} & \multicolumn{2}{c}{$-5.9\pm8.7$} & \multicolumn{2}{c}{$-8.6\pm8.9$} & \multicolumn{2}{c}{$-15.4\pm9.4$} \\
\cmidrule(lr){3-4}\cmidrule(lr){5-6}\cmidrule(lr){7-8}
 & & MAE & inl@1\% & MAE & inl@1\% & MAE & inl@1\% \\
\midrule
L0 & 单像素（Coates+MF） & 0.716 & 39.8 & 0.860 & 33.5 & 1.727 & 13.2 \\
L1 & 不对齐 + 等权平均 & 0.084 & 62.6 & 0.099 & 62.0 & 0.199 & 57.7 \\
L2 & 不对齐 + 信号强度权重 & 0.082 & 62.7 & 0.098 & 62.1 & 0.196 & 58.0 \\
L3 & 无先验 + 全轴匹配对齐 & 1.571 & 29.5 & 1.653 & 25.4 & 2.055 & 12.2 \\
L4 & 先验整数对齐 + 几何门 & 0.064 & 64.3 & 0.080 & 63.3 & 0.194 & 57.1 \\
L5 & L4 + 匹配门 + 信号强度 & 0.061 & 65.1 & 0.077 & 64.0 & 0.190 & 57.7 \\
L6 & L5 + 窗搜索重对齐 & \textbf{0.058} & \textbf{64.3} & \textbf{0.074} & \textbf{63.9} & \textbf{0.186} & \textbf{57.6} \\
\bottomrule
\end{tabular}
\end{table}

\begin{table}[!t]
\centering
\caption{46×62 分辨率下的消融实验结果（464 个场景全集；MAE 单位为 m，inl@1\% 为相对误差小于 1\% 的像素占比）}
\label{tab:ablation46}\scriptsize\setlength{\tabcolsep}{2.5pt}
\begin{tabular}{clcccccc}
\toprule
\multirow{2}{*}{级别} & \multirow{2}{*}{配置} & \multicolumn{2}{c}{$-5.7\pm8.4$} & \multicolumn{2}{c}{$-8.4\pm8.6$} & \multicolumn{2}{c}{$-15.2\pm9.0$} \\
\cmidrule(lr){3-4}\cmidrule(lr){5-6}\cmidrule(lr){7-8}
 & & MAE & inl@1\% & MAE & inl@1\% & MAE & inl@1\% \\
\midrule
L0 & 单像素（Coates+MF） & 0.665 & 37.0 & 0.814 & 30.9 & 1.695 & 12.5 \\
L1 & 不对齐 + 等权平均 & 0.323 & 22.7 & 0.349 & 22.3 & 0.522 & 19.7 \\
L2 & 不对齐 + 信号强度权重 & 0.317 & 23.0 & 0.342 & 22.5 & 0.519 & 20.0 \\
L3 & 无先验 + 全轴匹配对齐 & 1.470 & 27.3 & 1.533 & 23.2 & 2.031 & 10.0 \\
L4 & 先验整数对齐 + 几何门 & 0.214 & 29.3 & 0.235 & 30.1 & 0.460 & 24.1 \\
L5 & L4 + 匹配门 + 信号强度 & 0.197 & 32.4 & 0.221 & 32.5 & 0.451 & 25.2 \\
L6 & L5 + 窗搜索重对齐 & \textbf{0.184} & \textbf{40.1} & \textbf{0.207} & \textbf{38.2} & \textbf{0.427} & \textbf{27.5} \\
\bottomrule
\end{tabular}
\end{table}

结果呈现三个一致的规律。其一，无先验的数据驱动对齐（L3）在两个分辨率下均系统性崩溃（MAE $1.5\sim2.1\,\mathrm{m}$，甚至劣于不融合的 L0）：归一化匹配得分在背景主导的直方图上频繁锁定背景巧合峰，且跨邻居的错位独立发生、无法通过平均抵消。将搜索中心约束为先验整数对齐（L4）后误差立即恢复一个数量级以上，验证了 \ref{sec:alignment} 节“单目先验提供搜索范围、数据决定最终对齐”的分工必要性；对照式~\eqref{eq:delta}，此时起作用的仅是单目相对位移的几何趋势，其全局尺度误差不参与。其二，几何门（L4）与匹配门、信号强度（L5）的逐级引入在两个分辨率下均带来单调改善，说明几何门、匹配门与幅度响应的贡献可分离且方向一致。其三，窗搜索（L6）在统一 $K=5$ 配置下于两个分辨率、三个照度档全部取得最优 MAE 与中位数误差，相对 L5 的净增益为 $3\%\sim7\%$；同时 inl@1\% 不降，表明小窗重对齐在修正邻居残余错位的同时未引入背景巧合峰。

值得指出的是 MAE 与 inl@1\% 在 L0$\to$L1 处的背离：46×62 上等权平均使 MAE 从 $0.665$ 降至 $0.323$，inl@1\% 却从 $37.0\%$ 跌至 $22.7\%$。这是误差分布重组的结果——单像素估计呈双峰分布（约 37\% 的像素亚厘米级准确，其余误差约 $1\,\mathrm{m}$ 量级），等权平均将高信噪比像素与跨边缘、饥饿邻居无差别混合，以头部精度换取尾部均值；后续几何门与匹配门逐级将该头部精度找回，至 L6 时 inl@1\% 反超 L0（$40.1\%$ vs $37.0\%$），实现两端兼顾。这一现象也说明单靠 MAE 会高估朴素融合的价值。


\subsection{对比实验}
\label{sec:comparison}

将本文方法与四类基线在完全相同的 464 个场景、同一真值与同一评价脚本下比较：（1）Argmax，原始直方图直接取峰（无 Coates 校正与匹配滤波）；（2）Raw+MF，原始直方图经匹配滤波后取峰；（3）Coates+MF，即消融实验中的 L0；（4）CASPI~\cite{Lee2023}，使用官方实现与原论文参数；（5）本文方法（L6，$K=5$）。结果见表~\ref{tab:comp230}与表~\ref{tab:comp46}。

\begin{table}[t]
\centering
\caption{230×310 分辨率下与基线方法的定量对比（464 个场景全集；上半部分为 MAE，单位 m；下半部分为 inl@1\%，单位 \%；列组标题为全像幅平均检测 SNR，mean$\pm\mathrm{std}$，单位 dB）}
\label{tab:comp230}\scriptsize\setlength{\tabcolsep}{6.5pt}
\begin{tabular}{l cc cc cc}
\toprule
 & \multicolumn{2}{c}{$-5.7 \pm 8.4$~dB} & \multicolumn{2}{c}{$-8.4 \pm 8.6$~dB} & \multicolumn{2}{c}{$-15.2 \pm 9.0$~dB} \\
\cmidrule(lr){2-3} \cmidrule(lr){4-5} \cmidrule(lr){6-7}
Method & MAE & inl@1\% & MAE & inl@1\% & MAE & inl@1\% \\
\midrule
Argmax    & 1.645 & 18.28 & 1.724 & 14.85 & 1.976 & 7.19 \\
MF        & 0.404 & 39.56 & 0.553 & 33.14 & 1.048 & 18.08 \\
Coates+MF & 0.665 & 37.04 & 0.814 & 30.94 & 1.694 & 12.48 \\
CASPI     & 0.207 & 38.07 & 0.282 & 35.02 & 0.648 & 25.43 \\
Ours      & \textbf{0.184} & \textbf{40.12} & \textbf{0.207} & \textbf{38.22} & \textbf{0.427} & \textbf{27.46} \\
\bottomrule
\end{tabular}
\end{table}

\begin{table}[t]
\centering
\caption{46×62 分辨率下与基线方法的定量对比（464 个场景全集；上半部分为 MAE，单位 m；下半部分为 inl@1\%，单位 \%；列组标题为全像幅平均检测 SNR，mean$\pm\mathrm{std}$，单位 dB）}
\label{tab:comp46}\scriptsize\setlength{\tabcolsep}{6.5pt}
\begin{tabular}{l cc cc cc}
\toprule
 & \multicolumn{2}{c}{$-5.9 \pm 8.7$~dB} & \multicolumn{2}{c}{$-8.6 \pm 8.9$~dB} & \multicolumn{2}{c}{$-15.4 \pm 9.4$~dB} \\
\cmidrule(lr){2-3} \cmidrule(lr){4-5} \cmidrule(lr){6-7}
Method & MAE & inl@1\% & MAE & inl@1\% & MAE & inl@1\% \\
\midrule
Argmax    & 1.645 & 18.25 & 1.728 & 14.76 & 1.980 & 7.17 \\
MF        & 0.415 & 43.14 & 0.559 & 36.26 & 1.073 & 20.20 \\
Coates+MF & 0.716 & 39.84 & 0.860 & 33.48 & 1.727 & 13.23 \\
CASPI     & \textbf{0.097} & 53.39 & \textbf{0.140} & 51.65 & \textbf{0.375} & 43.63 \\
Ours      & \textbf{0.058} & \textbf{64.28} & \textbf{0.074} & \textbf{63.95} & \textbf{0.186} & \textbf{57.60} \\
\bottomrule
\end{tabular}
\end{table}

对比给出三点观察。首先，传统预处理的选择比直觉敏感：Coates 校正作为反卷积类操作，在光子饥饿条件下将稀疏计数中的噪声放大为高频伪峰，使 Coates+MF 系统性劣于不校正的 Raw+MF（MAE 劣化 $47\%\sim73\%$），匹配滤波在此承担的实为抑制该校正副作用的角色；该差异在 argmax 级同样存在（不校正直方图直接取峰优于 Coates 后取峰 2 倍以上），说明这是校正本身的统计性质而非某一实现的问题。其次，CASPI 在 46×62 上是强基线（较单像素基线改善 $55\%\sim69\%$），但在 230×310 上相对自身退化、与本文方法的差距拉大至 $40\%\sim51\%$：其非局部校正依赖像素间统计相似性，原生分辨率下单像素光子数下降一个量级后相似性证据随之减弱，而本文方法因单目先验提供了与光子统计无关的几何约束，受分辨率与照度的影响显著更小。最后，在全部 12 个工作点中本文方法均取得最优 MAE 与 inl@1\%，且尾部指标的领先幅度大于均值指标（230×310 下 inl@1\% 较 CASPI 高 $6\sim14$ 个百分点），与消融实验中融合权重对不可融合邻居的抑制作用一致。效率方面，在相同的 RTX 2080 Ti 上，CASPI 处理单帧需 89.6 s（230×310），本文完整融合流水线（含单目推理）约 55 s，且该耗时同时完成了全部八个消融变体的计算。

   % 实验章节（独立文件，便于单独修改）

\section{结论}
本模板覆盖了中文双栏论文最常用的排版元素，替换文字与图表即可开始写作。引用示例：单引\cite{spad2026}，多引\cite{nyudepth,lshort}。

% ================= 参考文献 =================
% 使用 BibTeX：引用处写 \cite{ref1}，条目放在 refs.bib
% Overleaf 会自动跑 bibtex；IEEEtran.bst 为 IEEE 官方样式
\bibliographystyle{IEEEtran}
\bibliography{refs}

\end{document}
